\scanpage{38}
(qui implique déjà $g_0=\rho g_\infty\ \deleted{\unsure{g_0}}$ dans (17)).

Les $\tilde u$ en question correspondent donc aux systèmes $(p,\sigma,g_0)\in\hat{\mathbb Z}^*\times\unsure{\mathfrak S_3^+|\mathcal S_3^+}\times\hat\pi$ tels que l'on ait la relation qui correspond à (\unsure{2})
\[
(18)\qquad\operatorname{int}(\rho^2(g_0))(l_{\sigma\infty}^p)\cdot\operatorname{int}(\rho(g_0))(l_{\sigma1}^p)\cdot\operatorname{int}(g_0)(l_{\sigma(0)}^p)=1
\]
(et que les trois facteurs engendrent $\hat\pi$\ldots)

b) Pour que $\tilde u$ commute à $\tilde\sigma_0$, il f. et s. que $\sigma$ y commute à $\dot\sigma_0$ (i.e. $\sigma\in\{1,\sigma_0\}$), et qu'on ait
\[
g_0=\sigma_0(g_0),\quad g_\infty=\sigma_0(g_1),\quad\inserted{\text{(donc)}}\ g_1=\sigma_0(g_\infty)
\]
\margin{moins intéressant pour la suite.}

\deleted{\unsure{On remarque d'abord}}

\deleted{\unsure{exprimons} que $\tilde u$ commute}

À vrai dire, nous nous intéressons plutôt à la commutation dans $\widehat{\widehat\Gamma}$, i.e. la commutation mod autom. intérieurs.
\mspage{37}
